\documentclass[10pt,a4paper]{amsart}
\usepackage[margin=19mm]{geometry}
\usepackage{amsmath,amssymb,mathtools}
\usepackage[scr=rsfs,cal=euler]{mathalfa}
\usepackage{xfrac}
\usepackage[unicode,hidelinks,bookmarks=false]{hyperref}
\newcommand{\ie}{i.e.\ }
\newcommand{\cf}{cf.\ }
\newcommand{\resp}{respectively\ }
\newcommand{\Subsection}{\subsection}
\providecommand{\nobreakdash}{\nobreak\hbox{-}}
\setlength{\emergencystretch}{2em}
\multlinegap0pt
\usepackage{amssymb}
\usepackage{mathtools}
\usepackage{bm}
\usepackage[shortlabels]{enumitem}
\newtheorem{theorem}{Theorem}
\title{\textbf{Source-Driven Transitions in Axisymmetric Euler Flow}}
\author{Rishad Shahmurov}
\date{Source: arXiv:2605.04526v2 (CC BY 4.0)}
\begin{document}
\maketitle

\noindent\emph{Source and licence.} \textbf{Source-Driven Transitions in Axisymmetric Euler Flow}. Version arXiv:2605.04526v2 (CC BY 4.0), distributed by arXiv under the Creative Commons Attribution 4.0 International licence, http://creativecommons.org/licenses/by/4.0/, as recorded in the arXiv licence metadata for this exact version and shown on the abstract page https://arxiv.org/abs/2605.04526v2. Full licence notice: \texttt{licenses/arxiv-2605.04526-CC-BY-4.0.txt}.\par\smallskip

\noindent\textit{Editorial scope.} Counted unit: the article's theorem thm:reference-derivatives (source0.tex lines 728--747). The article's complete original proof of it is delivered with it (source0.tex lines 748--781, one \texttt{proof} environment of 34 lines). No mathematics is written, repaired or re-derived: the statement and the proof are the article's own lines, in the article's own notation, and no retained line is substituted or re-flowed.  No further source-local context is needed: every cross-reference of every retained line resolves inside the two delivered spans.  Nothing was omitted from any retained span, so the package carries no omission record.  No retained line contains a citation, so the wrapper carries no bibliography and no bibliography entry.  No retained line contains a figure environment, an image-inclusion command or drawing code, so no image is delivered and the package declares no asset dependency.  Editorial formatting only, in addition: an AMS \texttt{amsart} preamble that restates the article's own declarations and macros used by the retained lines, namely the environments \texttt{theorem} on separate counters, together with the standard packages those lines need.  The retained environments print in this document as Theorem 1 in order of appearance, whereas the article prints the same items with the numbering of its own full document.  The complete native source of the article as distributed in the arXiv e-print for this exact version is bundled unchanged as \texttt{sources/199867/source0.tex}.
\begin{theorem}\label{thm:reference-derivatives}
For every $0\le s\le T$ the finite reference satisfies
\[
\begin{array}{c|c|c|c}
\text{region}&\text{quantity}&\text{certified bound}&\text{comparison cap}\\ \hline
B_{0.6}&|\nabla\bar\omega|&2.798&4\\
B_{0.6}&|D^2\bar\omega|_F&6.874&16\\
B_{0.6}&|\nabla\bar\theta|&44.027&64\\
B_{0.6}&|D^2\bar\theta|_F&19.144&64\\
B_{0.6}&|D^3\bar\theta|_F&128.642&1024\\
B_{0.9}&|\nabla\bar\omega|&4.475&8\\
B_{0.9}&|D^2\bar\omega|_F&11.549&32\\
B_{0.9}&|\nabla\bar\theta|&49.105&64\\
B_{0.9}&|D^2\bar\theta|_F&35.332&64\\
B_{0.9}&|D^3\bar\theta|_F&1219.173&2048\\
B_1&|\nabla\bar\omega|&5.393&8.
\end{array}
\]
The underlying certificate contains forty-four componentwise inequalities; the table lists the grouped maxima that dominate them.
\end{theorem}

\begin{proof}
Write one Bernstein coefficient field as a finite Fourier sum
\[
f(x,y)=\sum_{(k,l)\in\Lambda} c_{k,l}e^{i\alpha(kx+ly)},
\]
where the coefficients $c_{k,l}$ are stored dyadic numbers.  For a multi-index $\gamma$,
\[
D^\gamma f(x,y)=\sum_{(k,l)\in\Lambda}
(i\alpha k)^{\gamma_1}(i\alpha l)^{\gamma_2}c_{k,l}e^{i\alpha(kx+ly)}.
\]
Thus every derivative is again a finite Fourier sum with exact algebraic frequency multipliers.  The numerical proof converts the dyadic coefficients exactly, encloses the trigonometric factors, and performs all additions with outward rounding.  Hence the interval at a grid node contains the exact value of the corresponding derivative of the finite reference.

Let $x_g$ be a certified grid node with $|x-x_g|\le\delta_g$.  Taylor's formula in integral form gives
\[
D^\gamma f(x)=D^\gamma f(x_g)
+\int_0^1D^{\gamma+1}f(x_g+t(x-x_g))[x-x_g]\,dt.
\]
The computation includes the derivative order needed to bound the integrand.  Consequently
\[
|D^\gamma f(x)|
\le |D^\gamma f(x_g)|+\delta_g
\sup_{[x_g,x]}|D^{\gamma+1}f|.
\]
Applying this componentwise and using the Euclidean or Frobenius norm gives the continuum spatial bounds in the table.  Finally, if
\[
\bar f(s)=\sum_{j=0}^3B_j(s)f_j,
\qquad B_j(s)\ge0,\qquad\sum_jB_j(s)=1,
\]
then convexity of the norm yields
\[
|D^\gamma\bar f(s)|\le\max_j|D^\gamma f_j|.
\]
This proves uniformity in time and completes the argument.
\end{proof}

\end{document}
